\documentclass[11pt, a4paper]{article}

% --- UNIVERSAL PREAMBLE BLOCK ---
\usepackage[a4paper, top=2.5cm, bottom=2.5cm, left=2cm, right=2cm]{geometry}
\usepackage{fontspec}
\usepackage[spanish, bidi=basic, provide=*]{babel}

\babelprovide[import, onchar=ids fonts]{spanish}
\babelprovide[import, onchar=ids fonts]{english}

% Configuración de la fuente Noto Sans
\babelfont{rm}{Noto Sans}

% Paquetes técnicos
\usepackage{amsmath}
\usepackage{booktabs}
\usepackage{graphicx}
\usepackage{listings}
\usepackage{xcolor}
\usepackage[siunitx, american]{circuitikz}
\usepackage{tikz}
\usetikzlibrary{shapes.geometric, arrows, positioning, calc}
\usepackage{enumitem}
\usepackage{hyperref}
\usepackage{array}
\usepackage{caption}

% --- ESTILO DE CÓDIGO ---
\lstset{
    language=Python,
    basicstyle=\ttfamily\tiny,
    keywordstyle=\color{blue},
    stringstyle=\color{red},
    commentstyle=\color{green!60!black},
    breaklines=true,
    numbers=left,
    numberstyle=\tiny\color{gray},
    frame=single,
    backgroundcolor=\color{gray!5}
}

% --- ESTILOS TIKZ ---
\tikzset{
    startstop/.style={rectangle, rounded corners, minimum width=3cm, minimum height=1cm, text centered, draw=black, fill=red!20, font=\small},
    process/.style={rectangle, minimum width=4.5cm, minimum height=1.5cm, text width=4.2cm, text centered, draw=black, fill=orange!15, font=\small},
    decision/.style={diamond, aspect=2.2, minimum width=3.8cm, minimum height=1cm, text centered, draw=black, fill=green!15, font=\small},
    arrow/.style={thick,->,>=stealth}
}

\begin{document}

% --- PORTADA ---
\begin{titlepage}
    \centering
    {\LARGE Universidad Autónoma de Nuevo León}\\[0.5cm]
    {\Large Facultad de Ingeniería Mecánica y Eléctrica}\\[2cm]
    
    {\Large Laboratorio de Técnicas de Medida en Aeronáutica}\\[1cm]
    
    {\Huge Práctica 8}\\[0.5cm]
    {\Large Análisis de Productos de Combustión y Eficiencia Energética}\\[2cm]
    
    Elaborado para el ciclo escolar:\\
    2026\\[2cm]
    
    \vfill
    
    Resumen Ejecutivo\\
    Esta práctica aborda la caracterización de la eficiencia de combustión en motores de reacción aeronáuticos mediante el análisis químico de los gases de escape. El estudiante procesará conjuntos de datos de concentración de especies ($CO_2$, $CO$, $O_2$, $NO_x$, $UHC$) bajo diversas condiciones operativas. El enfoque analítico se centra en el cálculo de la relación aire/combustible real, el coeficiente de exceso de aire ($\lambda$) y la evaluación del impacto ambiental y operativo del proceso de oxidación del queroseno.
\end{titlepage}

\tableofcontents
\newpage

% --- NOMENCLATURA ---
\section*{Nomenclatura}
\begin{description}
    \item[$AFR$] Relación Aire/Combustible (Air-Fuel Ratio) [kg$_{aire}$/kg$_{comb}$]
    \item[$\lambda$] Coeficiente de exceso de aire (adimensional)
    \item[$\eta_c$] Eficiencia de combustión [\%]
    \item[$ppm$] Partes por millón (concentración volumétrica)
    \item[$M_{fuel}$] Masa molar del combustible [$g/mol$]
    \item[$Y_i$] Fracción molar de la especie $i$
    \item[$UHC$] Hidrocarburos no quemados (Unburned Hydrocarbons)
    \item[$EGT$] Temperatura de gases de escape (Exhaust Gas Temperature)
\end{description}

\section{Introducción}
El análisis de los productos de combustión es una herramienta fundamental para diagnosticar la salud de un motor de reacción y asegurar el cumplimiento de normativas ambientales internacionales (OACI). La composición química del escape revela no solo la eficiencia térmica del ciclo, sino también problemas latentes en el sistema de inyección o en la integridad de la cámara de combustión. En esta práctica, el estudiante actuará como analista de rendimiento, convirtiendo datos de concentración volumétrica en parámetros de decisión ingenieril.

\section{Objetivos de Aprendizaje}
\begin{itemize}[label=-]
    \item Identificar las especies químicas críticas resultantes de la combustión del queroseno ($C_{12}H_{23}$).
    \item Realizar balances de masa de carbono y nitrógeno para determinar la relación aire/combustible real ($AFR_{actual}$).
    \item Evaluar el exceso de aire ($\lambda$) y su correlación con la estabilidad de la llama y las emisiones de $NO_x$.
    \item Estimar la eficiencia de combustión ($\eta_c$) cuantificando las pérdidas por energía química no liberada ($CO$ y $UHC$).
    \item Discutir las implicaciones del régimen de operación del motor en la firma de emisiones.
\end{itemize}

\section{Fundamento Teórico}

\subsection{Química de la Combustión Aeronáutica}
El queroseno aeronáutico se modela típicamente como una molécula promedio de $C_{12}H_{23}$. La combustión completa e ideal sigue la reacción estequiométrica:
\begin{equation}
    C_{12}H_{23} + 17.75(O_2 + 3.76N_2) \rightarrow 12CO_2 + 11.5H_2O + 66.74N_2
\end{equation}
A partir de esta relación, se define el $AFR_{stoich} \approx 14.57$. En condiciones reales, la presencia de $CO$ y hidrocarburos sin quemar indica una combustión incompleta, lo cual reduce la energía cinética disponible para la propulsión.

\subsection{Instrumentación de Muestreo de Gases}
El sistema de medida requiere una sonda de muestreo resistente a altas temperaturas, un sistema de enfriamiento y condensación (trampa de agua) y analizadores específicos como NDIR (Infrarrojo No Dispersivo) para $CO/CO_2$ y FID (Detección por Ionización de Llama) para hidrocarburos.

\begin{figure}[htbp]
    \centering
    \begin{tikzpicture}[scale=0.9]
        % Cámara de combustión
        \draw[thick] (0,0) rectangle (3,2);
        \node at (1.5,1) {Motor};
        
        % Sonda
        \draw[line width=1.5pt] (3,1) -- (5,1);
        \node at (4,1.3) [font=\tiny] {Sonda de Gases};
        
        % Línea de transporte
        \draw[->, thick] (5,1) -- (6,1) -- (6,0) -- (7,0);
        
        % Acondicionador
        \draw[fill=blue!10] (7,-0.5) rectangle (9,0.5);
        \node at (8,0) [font=\tiny, align=center] {Trampa de\\Agua (Secado)};
        
        % Analizadores
        \draw[->, thick] (9,0) -- (10,0);
        \draw[fill=gray!20] (10,-1) rectangle (13,1);
        \node at (11.5,0) [font=\tiny, align=center] {Analizadores\\$CO_2, CO, O_2, NO_x, UHC$};
        
        \node at (6.5, -1.5) [font=\scriptsize, text width=10cm, align=center] {Esquema de la cadena de medida: la muestra debe ser filtrada y secada para el análisis en base seca.};
    \end{tikzpicture}
    \caption{Sistema de adquisición de datos para productos de combustión.}
\end{figure}

\section{Metodología de Análisis de Datos}

\subsection{Algoritmo de Procesamiento Térmico-Químico}
El estudiante debe seguir el flujo lógico para transformar las concentraciones medidas en parámetros de eficiencia:

\begin{figure}[htbp]
    \centering
    \begin{tikzpicture}[node distance=1.7cm, auto]
        \node (input) [startstop] {Concentraciones ($CO_2, O_2, CO, NO_x, UHC$)};
        \node (n2) [process, below=of input] {Cálculo de $N_2$ (Balance por diferencia en base seca)};
        \node (afr) [process, below=of n2] {Determinación de $AFR_{actual}$ (Balance C-N)};
        \node (lambda) [process, below=of afr] {Cálculo de $\lambda = AFR_{actual} / AFR_{stoich}$};
        \node (eff) [process, below=of lambda] {Cálculo de $\eta_c$ (Pérdidas químicas)};
        \node (rep) [startstop, right=of eff, xshift=2cm] {Informe AIAA};
        
        \draw [arrow] (input) -- (n2);
        \draw [arrow] (n2) -- (afr);
        \draw [arrow] (afr) -- (lambda);
        \draw [arrow] (lambda) -- (eff);
        \draw [arrow] (eff) -- (rep);
    \end{tikzpicture}
    \caption{Flujo metodológico para la evaluación del proceso de combustión.}
\end{figure}

\newpage
\section{Casos de Estudio y Datasets}

\subsection{Datos de Emisiones por Régimen de Potencia}
Se proporcionan las mediciones de un motor turbofán operando bajo el ciclo LTO (Landing and Take-Off).

\begin{table}[htbp]
    \centering
    \begin{tabular}{lcccccc}
        \toprule
        Condición & Potencia [\%] & $CO_2$ [\%] & $O_2$ [\%] & $CO$ [ppm] & $NO_x$ [ppm] & $UHC$ [ppm] \\
        \midrule
        Ralentí (Idle) & 25 & 7.5 & 11.5 & 2500 & 60 & 350 \\
        Aproximación & 40 & 9.0 & 9.0 & 1000 & 150 & 150 \\
        Crucero & 75 & 12.0 & 5.0 & 400 & 700 & 40 \\
        Despegue & 100 & 13.8 & 1.8 & 700 & 1800 & 70 \\
        \bottomrule
    \end{tabular}
    \caption{Concentraciones medidas en los gases de escape (base seca).}
\end{table}

\section{Actividades a Realizar}
\begin{enumerate}
    \item Cálculo de Nitrógeno: Determine el porcentaje de $N_2$ para cada condición asumiendo que el resto de los gases secos no medidos son despreciables.
    \item Análisis de Mezcla: Calcule el $AFR_{actual}$ utilizando la masa molar del aire ($28.97\,g/mol$) y del queroseno ($167.316\,g/mol$).
    \item Evaluación de Eficiencia: Aplique la fórmula de eficiencia de combustión considerando el poder calorífico relativo del monóxido de carbono.
    \item Modelado de Emisiones: Utilice el script de Python para graficar las tendencias de contaminantes frente a la potencia del motor.
\end{enumerate}

\section{Análisis de Resultados y Graficación}
En el reporte técnico se deben discutir las siguientes situaciones de estudio:
\begin{enumerate}
    \item Relación Potencia-Emisiones: Explique por qué el $CO$ y los $UHC$ son máximos en ralentí, mientras que el $NO_x$ alcanza su pico en despegue. Relacione esto con la temperatura de la llama.
    \item Sensibilidad de la Eficiencia: Grafique $\eta_c$ vs. Potencia. Identifique el régimen de máxima eficiencia y discuta el costo económico de operar a baja potencia.
    \item Análisis de Exceso de Aire: Discuta los valores de $\lambda$ obtenidos. ¿Por qué los motores de turbina operan con mezclas mucho más pobres que los motores de pistón?
\end{enumerate}

\section{Preguntas de Discusión Electrónica y de Combustión}
\begin{enumerate}
    \item ¿Cómo influye el tiempo de respuesta de los sensores de gases en la capacidad de detectar inestabilidades de combustión (Rumble o Screech) en tiempo real?
    \item Explique el principio físico de un sensor de Oxígeno paramagnético y por qué es preferido sobre sensores electroquímicos en bancos de prueba de alta precisión.
    \item En la instrumentación de medida, ¿por qué es crítico eliminar el vapor de agua de la muestra antes de pasar por los analizadores de infrarrojo ($CO_2$)?
    \item Si se detecta un aumento súbito de $UHC$ en crucero, ¿qué falla electrónica o mecánica en el sistema de combustible podría ser la causa raíz?
\end{enumerate}

\section{Lineamientos del Reporte (AIAA)}
El reporte debe presentar las tablas de resultados calculados ($AFR, \lambda, \eta_c$) para cada régimen. Se debe incluir la comparativa gráfica de emisiones y una sección dedicada al diagnóstico del motor basada en las tendencias observadas, vinculando los resultados con la eficiencia energética global de la aeronave.

\newpage
\appendix
\section{Apéndice: Script de Procesamiento (\texttt{analisis\_combustion.py})}
Script para automatizar los cálculos de balance de masa y generación de perfiles de emisiones.

\begin{lstlisting}[language=Python]
import pandas as pd
import numpy as np

# Constantes del Combustible
M_FUEL = 167.316 
M_AIR = 28.97
AFR_STOICH = 14.57
C_ATOMS = 12

def calcular_parametros(co2_pct, o2_pct, co_ppm, uhc_ppm):
    # 1. Calculo de N2 por diferencia
    n2_pct = 100 - co2_pct - o2_pct - (co_ppm/10000) - (uhc_ppm/10000)
    
    # 2. Fracciones molares
    y_co2 = co2_pct / 100
    y_co = co_ppm / 1e6
    y_uhc = uhc_ppm / 1e6
    
    # 3. AFR Actual (Balance C-N)
    den = y_co2 + y_co + y_uhc
    afr = (M_AIR / M_FUEL) * (C_ATOMS * (n2_pct / 79.05)) / den
    
    # 4. Eficiencia de Combustion (Aprox)
    eta = (1 - (0.427 * co_ppm + uhc_ppm) / (co2_pct * 10000 + co_ppm + uhc_ppm)) * 100
    
    return afr, eta

# El script debe iterar sobre la tabla de datos y generar graficas
\end{lstlisting}

\section{Referencias}
\begin{enumerate}
    \item Mattingly, J. D. (2006). Elements of Propulsion: Gas Turbines and Rockets. AIAA Education Series.
    \item Lefebvre, A. H., \& Ballal, D. R. (2010). Gas Turbine Combustion: Alternative Fuels and Emissions.
    \item Heywood, J. B. (1988). Internal Combustion Engine Fundamentals. McGraw-Hill.
\end{enumerate}

\end{document}